\documentclass[12pt,a4paper]{article}

\usepackage[margin=1in]{geometry}
\usepackage{enumitem}
\usepackage{hyperref}
\usepackage{titlesec}

\titleformat{\section}{\large\bfseries}{\thesection.}{0.5em}{}
\titleformat{\subsection}{\normalsize\bfseries}{\thesubsection}{0.5em}{}

\title{
  \textbf{CampusLink}\\
  \large A Goal-Matching \& Resource-Sharing Platform for College Campuses\\
  \vspace{0.5cm}
  \normalsize Project Proposal for UCS503
}
\author{
  Govind, CSED -- Roll No: TBD, Email: TBD\\
  Teammate 1, CSED -- Roll No: TBD, Email: TBD\\
  Teammate 2, CSED -- Roll No: TBD, Email: TBD
}
\date{August 2026}

\begin{document}

\maketitle

\section{Higher-Order Goal}

The broader aim of this project is to make it easier for students on a large campus to find people to do things with, whether it is a spontaneous plan, a recurring need like a gym partner, or sharing class notes. CampusLink tries to surface students who want the same thing at the same time and gives the interaction enough structure to feel safer than an informal group message.

\section{Problem Statement}

On a residential campus with thousands of students, everyday coordination still happens through fragmented informal channels such as batch WhatsApp groups, hostel groups, and word of mouth.

\begin{itemize}[leftmargin=*]
  \item \textbf{Limited reach:} students miss compatible plans because they do not know the right people.
  \item \textbf{No structure for stranger interaction:} one-on-one or late-night plans need more trust signals.
  \item \textbf{Inefficient resource sharing:} notes and academic resources are scattered across groups.
  \item \textbf{No accountability:} repeated no-shows and unreliable participation are hard to track.
\end{itemize}

\section{Proposed Solution}

CampusLink is a companion-matching platform built around a simple loop: a student posts a goal, the goal appears on the public campus feed, interested students request to join, and once enough people are accepted, a lobby forms. The same post, request, and fulfilment loop also supports notes and resource sharing through private chat.

\subsection{Core Workflow}

\begin{enumerate}[leftmargin=*]
  \item \textbf{Post:} create a goal with category, description, date/time, and headcount.
  \item \textbf{Discover:} browse and filter goals by category and date.
  \item \textbf{Request:} request to join, with poster approval or auto-accept.
  \item \textbf{Coordinate and close:} chat in-app, meet, confirm completion, and leave feedback.
\end{enumerate}

\subsection{Key Features}

\textbf{MVP features}
\begin{itemize}[leftmargin=*]
  \item College email verification.
  \item Goal posting with category, description, date/time, and headcount.
  \item Public feed with filters.
  \item Request-to-join flow.
  \item Basic profiles.
  \item Lobby chat and private resource-request chat.
\end{itemize}

\textbf{Trust and safety features}
\begin{itemize}[leftmargin=*]
  \item Report and block controls.
  \item Gender-preference filters for specific goals.
  \item Anonymous posting on the public feed before acceptance.
  \item Visible reliability score.
  \item No-show flagging.
\end{itemize}

\section{Solution Approach}

CampusLink will be implemented as a standard 3-tier web application. Matching will remain rule-based so the project stays focused on workflow correctness, system design, and user safety.

\section{Architecture}

\begin{itemize}[leftmargin=*]
  \item \textbf{Frontend:} React, Vite, Tailwind CSS, React Router, Zustand or Redux Toolkit.
  \item \textbf{Backend:} Node.js with Express or NestJS.
  \item \textbf{Database:} MongoDB or PostgreSQL.
  \item \textbf{Real-time layer:} Socket.IO.
  \item \textbf{Authentication:} JWT with college email OTP verification.
  \item \textbf{File storage:} Cloudinary or Firebase.
  \item \textbf{Deployment:} Vercel for frontend and Render/Railway for backend.
\end{itemize}

\section{Evaluation Criteria}

\textbf{Primary metric:} lobby fulfilment rate.

\textbf{Secondary metrics:}
\begin{itemize}[leftmargin=*]
  \item Time-to-first-request.
  \item No-show rate.
  \item Resource-request response rate.
  \item User-reported clarity and safety.
\end{itemize}

\section{Project Scope and Deliverables}

\subsection{Iteration 1 -- Core MVP}

Signup/login, profile pages, goal posting, feed filters, request-to-join workflow, lobby state management, and basic CI.

\subsection{Iteration 2 -- Trust, Safety, and Real-Time}

Socket.IO chat, report/block controls, gender-preference filter, anonymous posting, no-show flagging, match suggestions, ratings, and feedback.

\subsection{Iteration 3 -- Polish and Deploy}

Responsive UI, tests, deployment, README, and a short demo video.

\section{Risks and Mitigations}

\begin{itemize}[leftmargin=*]
  \item \textbf{Stranger meetup safety:} partially mitigated through reliability scores, report/block, and gender filters.
  \item \textbf{Cold-start problem:} seeded pilot goals will be used during testing.
  \item \textbf{Anonymity misuse:} admins can unmask identity for moderation after reports.
  \item \textbf{Low adoption:} the pilot will use a small test group within the course timeline.
\end{itemize}

\section{Summary}

CampusLink aims to solve an everyday campus problem: finding reliable people for shared goals beyond existing friend circles. The project focuses on the core workflow and trust layer needed to make stranger-to-stranger campus interaction practical within the UCS503P timeline.

\end{document}

